\begin{center}
\textbf{Une coulée comme Leon Marchand}
\end{center}

On appelle coulée la phase sous-marine qui suit le plongeon ou le virage d'un nageur.

On suppose que le déplacement du nageur se produit à une profondeur constante lors de la coulée. Pendant cette phase, le sportif se laisse glisser dans l'eau, sans nager. Le mouvement du nageur est supposé rectiligne.

La figure 1 ci-après représente l'évolution de la vitesse du nageur en fonction du temps.

\begin{center}
\psset{xunit=6cm,yunit=1.8cm,comma}
\begin{pspicture}(-0.2,-0.35)(2,4)
\psaxes[linewidth=1.25pt,Dx=0.2,Dy=0.5,labelFontSize=\scriptstyle](0,0)(0,0)(2,4)
\multido{\n=0.00+0.04}{51}{\psline[linewidth=0.1pt](\n,0)(\n,4)}
\multido{\n=0.0+0.2}{11}{\psline[linewidth=0.55pt](\n,0)(\n,4)}
\multido{\n=0.0+0.1}{41}{\psline[linewidth=0.1pt](0,\n)(2,\n)}
\multido{\n=0.0+0.5}{9}{\psline[linewidth=0.55pt](0,\n)(2,\n)}
\psplot[plotpoints=2000,linewidth=1.25pt,linecolor=red]{0}{1.78}{3.64 2.71828 0.6 x mul exp div}
\uput[d](1,-0.25){Temps (s)}\rput{90}(-0.2,2){Vitesse (m/s)}
\end{pspicture}
\medskip

\textbf{Figure 1 - Vitesse du nageur en fonction du temps
}
\end{center}

On détermine un modèle numérique à partir de  cette expérience. On suppose que la distance parcourue par le nageur durant la coulée, exprimée en mètre, en fonction du temps $t$, exprimée en seconde, est définie sur l'intervalle [0~;~2] par :

\[f(t) = 3,64\ln(1 + t).\]

\begin{enumerate}[start=5]
\item Calculer $f(0)$.

\item On admet que la vitesse du nageur, en m/s, est donnée, en fonction du temps en s, par $v(t) = \dfrac{3,64}{1 + t}$.

Montrer que l'accélération, correspondant à la dérivée de $v$ est $a(t) = - \dfrac{3,64}{(1 + t)^2}$.

\item Interpréter le signe de $a$ dans le contexte de l'exercice et vérifier la cohérence avec l'observation de la courbe représentative de la fonction $v$.
\end{enumerate}

\bigskip
